\section{Diseño e implementación secuencial}

La versión secuencial define la semántica de referencia y proporciona el tiempo base para calcular speedup.

El esqueleto implementado separa la interfaz de línea de comandos, la configuración y el registro en módulos comunes bajo \texttt{src/common}. El ejecutable \texttt{schelling\_seq} carga valores predeterminados, aplica el archivo de escenario, incorpora las sobrescrituras indicadas por argumentos y valida la configuración efectiva antes de iniciar el modelo. Esta organización permite que el ejecutable híbrido reutilice exactamente la misma interpretación de entradas.

\subsection{Estructuras de datos}

La grilla y los hogares se almacenan en arreglos contiguos separados. El identificador global de una celda coincide con su posición en el arreglo y sigue un recorrido por filas; por ello las coordenadas se derivan mediante división y resto sin almacenarlas en cada celda. Cada celda conserva tipo, zona, precio e identificador del hogar ocupante. El valor centinela $-1$ representa una vivienda vacía, evitando mantener un segundo indicador de ocupación que pudiera contradecirlo.

Los hogares conservan identificador, celda actual, subestrato, clase derivada, ingreso, bloqueo y satisfacción. Al crear o validar un estado se comprueba la correspondencia bidireccional entre hogar y celda, que una celda no residencial nunca esté ocupada y que cada hogar tenga exactamente una vivienda.

\subsection{Funciones y organización del código}

El módulo de modelo administra memoria, coordenadas, definición de celdas, ubicación de hogares e invariantes. El generador sintético construye el estado inicial a partir de parámetros y claves aleatorias estables. El módulo de vecindario precalcula una sola vez los desplazamientos de Chebyshev y sus pesos gaussianos. Sobre esas tablas calcula proporciones ponderadas y evalúa las tolerancias cargadas desde el escenario. \pendiente{completar las responsabilidades de simulación, precios, mudanzas y métricas.}

\subsection{Algoritmo}

Cada iteración evalúa satisfacción sobre la ocupación consolidada, actualiza solamente los precios de viviendas vacías y busca un destino para cada hogar insatisfecho no bloqueado. La búsqueda de referencia recorre todas las vacantes, descarta las inaccesibles o insatisfactorias y elige la de menor distancia, con identificador como desempate. Las solicitudes se almacenan sin modificar la grilla. Para cada vivienda se conserva el hogar de mayor ingreso; los empates usan una clave pseudoaleatoria estable y luego el identificador. Finalmente se liberan todos los orígenes ganadores y recién después se ocupan los destinos, por lo que una vivienda liberada no participa hasta el mes siguiente.

Para $H$ hogares, $V$ viviendas vacías y $K$ desplazamientos del vecindario, esta referencia tiene costo dominante $O(HVK)$ en el peor caso. Esta versión prioriza una semántica comprobable; el índice espacial posterior reduce la cantidad de destinos examinados. \pendiente{agregar el pseudocódigo definitivo después de incorporar el índice espacial.}

\subsection{Reproducibilidad y checkpoints}

El generador pseudoaleatorio no mantiene estado mutable. Una mezcla SplitMix64 recibe una clave compuesta por semilla, iteración, identificador de entidad, propósito y número de muestra. Los 53 bits superiores del resultado se convierten a un valor uniforme en el intervalo abierto $(0,1)$. Separar los propósitos de tipo de celda, ocupación, subestrato, ruido y desempate evita que agregar una muestra a una fase desplace las secuencias de las restantes. Existen vectores de prueba fijos para la función de mezcla. \pendiente{documentar la transformación normal y el formato final de checkpoint.}

En ejecuciones normales se genera un checkpoint cada 50 iteraciones y se conservan los dos últimos completos. La escritura usa archivos temporales y un manifiesto final para que un estado parcial no sea reiniciable. En los benchmarks los checkpoints se desactivan para evitar que la E/S altere la comparación de tiempos.

\subsection{Optimizaciones}

\pendiente{documentar optimizaciones únicamente después de medirlas y validarlas contra la referencia.}

\subsection{Convenciones y simplificaciones}

La grilla rectangular es una abstracción y no un mapa de Montevideo. Cada celda residencial representa una vivienda discreta, sin superficie física asociada. Los datos oficiales se conservan como cantidades y distribuciones agregadas por zona. Los índices comienzan en cero, las celdas se ordenan por filas y los bordes son cerrados: todo desplazamiento exterior se descarta.

\subsection{Dificultades de implementación}

\pendiente{describir problemas relevantes, diagnóstico y solución; omitir incidencias triviales que no aporten al análisis.}
